% Einheit 3 – Klammern auflösen (bank/terme/e3.jsonl, zusammenbau v1.0)
\einheitenkopf[e3][Klammern auflösen]{Klammern auflösen}
\setcounter{aufgabe}{14}


% test: terme-e3-k2-s8-v2
\begin{kbaufgabe}{Kannst du das schon? Dann weiter zu „Ausklammern“}
\begin{kbblock}
\kbzweispaltig[0.46]{\kbspalte{\kbanweisung{Löse die Klammern auf und fasse zusammen.}\kbteil{}{$5 \cdot (a + 3) - 2 \cdot (a - 1)$}{}\kbantwort{= \kbfeld}}}{\kbraum{2}}
\end{kbblock}
\end{kbaufgabe}


% erkennung: terme-e3-k1-s0-v1, terme-e3-k1-s0-v2, terme-e3-k1-s0-v3, terme-e3-k1-s0-v4
\begin{kbaufgabe}{Erkennen, was vor einer Klammer steht}
\begin{kbblock}
\kbteil{a)}{Was steht im Term $7 + (a - 2)$ direkt vor der Klammer?}{}
\kbantwort{\kbfeld}
\end{kbblock}
\begin{kbblock}
\kbteil{b)}{Was steht im Term $10 - (b + 4)$ direkt vor der Klammer?}{}
\kbantwort{\kbfeld}
\end{kbblock}
\begin{kbblock}
\kbteil{c)}{Was steht im Term $5 \cdot (x - 1)$ direkt vor der Klammer?}{}
\kbantwort{\kbfeld}
\end{kbblock}
\begin{kbblock}
\kbteil{d)}{Was steht im Term $2y - 6 \cdot (y + 1)$ direkt vor der Klammer?}{}
\kbantwort{\kbfeld}
\end{kbblock}
\end{kbaufgabe}


% leiter: terme-e3-k2-s1-v1, terme-e3-k2-s1-v2, terme-e3-k2-s1-v3, terme-e3-k2-s2-v1, terme-e3-k2-s3-v1, terme-e3-k2-s4-v1, terme-e3-k2-s5-v1, terme-e3-k2-s6-v1, terme-e3-k2-s7-v1
\begin{kbaufgabe}{Klammern auflösen}
\begin{kbblock}
\kbanweisung{Löse die Klammer auf.}
\lbpaar{\kbteil{a)}{$5a + (2b - 1) =$ \lbfeld}{}
}{\kbteil{b)}{$5a + (2b - 3) =$ \lbfeld}{}}
\end{kbblock}
\begin{kbblock}
\lbpaar{\kbteil{c)}{$5a + (2b - 4) =$ \lbfeld}{}
}{\item[]}
\end{kbblock}
\begin{kbblock}
\kbanweisung{Löse die Klammer mit der Tabelle auf.}
\kbteil{d)}{$5 \cdot (x + 3)$}{}
\kbzweispaltig[0.34]{\sachtabelle{ccc}{$\cdot$ & $x$ & $+3$}{$5$ & \leerzelle & \leerzelle}}{\lbantwortrechts[0.4cm]{= \kbfeld}}
\end{kbblock}
\begin{kbblock}
\kbteil{e)}{Rechne $30 - (12 + 8)$ auf zwei Wegen: erst die Klammer ausrechnen und dann abziehen; dann Glied für Glied abziehen. Vergleiche beide Ergebnisse.}{}
\item[]\kbraum{2}
\end{kbblock}
\begin{kbblock}
\kbanweisung{Löse die Klammer auf.}
\lbpaar{\kbteil{f)}{$5a - (2b + 4) =$ \lbfeld}{}
}{\kbteil{g)}{$7a - (2b - 5 + c) =$ \lbfeld}{}}
\end{kbblock}
\begin{kbblock}
\lbpaar{\kbteil{h)}{$-2 \cdot (x + 5) =$ \lbfeld}{}
}{\item[]}
\end{kbblock}
\begin{kbblock}
\kbzweispaltig[0.46]{\kbspalte{\kbanweisung{Löse die Klammer auf und fasse zusammen.}\kbteil{i)}{$5x + 2 \cdot (x + 6)$}{}\kbantwort{= \kbfeld}}}{\kbraum{2}}
\end{kbblock}
\end{kbaufgabe}


% pruefung: terme-e3-k2-s8-v1
\begin{kbaufgabe}{Auch zwei Klammern in einem Term auflösen und zusammenfassen}
\begin{kbblock}
\kbzweispaltig[0.46]{\kbspalte{\kbanweisung{Löse die Klammern auf und fasse zusammen.}\kbteil{}{$3 \cdot (2x + 5) - 2 \cdot (x - 4)$}{}\kbantwort{= \kbfeld}}}{\kbraum{2}}
\end{kbblock}
\end{kbaufgabe}


% pflicht: terme-e3-k3-s1-v1, terme-e3-k3-s2-v1, terme-e3-k3-s3-v1, terme-e3-k3-s4-v1
\begin{kbaufgabe}{Verstanden?}
\begin{kbblock}
\kbteil{a)}{Jan soll die Klammer in $14 - (a - 5)$ auflösen. Er rechnet so: \rechnung{14 - (a - 5) &= 14 - a - 5 &= 9 - a} Finde den Fehler und rechne richtig.}{}
\item[]\kbraum{2}
\end{kbblock}
\begin{kbblock}
\kbteil{b)}{Begründe, ob die Terme $5 - (x - 2)$ und $7 - x$ gleichwertig sind.}{}
\item[]\kbraum{3}
\end{kbblock}
\begin{kbblock}
\kbteil{c)}{Die Figur zeigt ein Rechteck. Die Maße sind in cm angegeben. Schreibe einen Term für den Flächeninhalt mit Klammer. Schreibe den Term dann ohne Klammer.}{}
\kbzweispaltig[0.48]{\viereck[seiten={x+2,5,x+2,5}]{(0,0)}{(6,0)}{(6,2.5)}{(0,2.5)}}{\lbantwortrechts[1.5cm]{\kbfeld}}
\end{kbblock}
\begin{kbblock}
\kbteil{d)}{Ein Kinoticket kostet 9 €. Eine Klasse mit 24 Schülern nimmt noch $x$ Begleiter mit. Schreibe die Gesamtkosten in Euro als Term mit Klammer. Löse die Klammer auf.}{}
\item[]\kbraum{2}
\end{kbblock}
\end{kbaufgabe}
